\section*{Overview}

Tree DP works because removing one edge splits the structure cleanly. That means a subtree often carries exactly the
local information needed for a recurrence, without the cycles that make general graph DP difficult.

\section*{When to Use It}

Use tree DP when:

\begin{itemize}
  \item the input is a tree or can be reduced to one,
  \item the answer decomposes naturally across child subtrees,
  \item a parent only needs a summary from each child, not the full internal structure.
\end{itemize}

\section*{Core Idea}

Root the tree, then define DP states per vertex. A typical pattern is:
\[
dp[v][state] = \text{best answer inside the subtree of } v \text{ under some local condition at } v.
\]

\section*{Key Insight}

The tree structure removes cyclic dependencies. Once the children are solved, their contributions can be merged into
the parent's state.

\section*{Operations / Main Technique}

\begin{itemize}
  \item choose a root,
  \item define a local state at each node,
  \item DFS from leaves upward,
  \item merge child contributions.
\end{itemize}

\paragraph{Worked example.}
For maximum independent set on a tree, use two states:
\begin{itemize}
  \item \(\texttt{take[v]}\): best if \(v\) is chosen,
  \item \(\texttt{skip[v]}\): best if \(v\) is not chosen.
\end{itemize}
If \(v\) is taken, no child may be taken. If \(v\) is skipped, each child chooses its better state.

\section*{Correctness Intuition}

Every edge is cut exactly once by the parent-child direction, so a subtree DP can summarize all decisions below a node
without ambiguity. The merge step is correct because child subtrees are disjoint except for the parent connection.

\section*{Complexity Analysis}

Many tree DP routines are \(O(n)\) or \(O(n \log n)\), depending on the merge cost per child.

\section*{Implementation}

The reference code implements maximum independent set on a tree, which is a clean example of the parent-child state
split and the merge logic.

\section*{Common Pitfalls}

\begin{itemize}
  \item Forgetting to exclude the parent during DFS.
  \item Choosing a state that is too weak and loses necessary information.
  \item Writing an \(O(n^2)\) merge when the state can be simplified.
\end{itemize}

\section*{Variants / Extensions}

\begin{itemize}
  \item Rerooting DP.
  \item DP on trees with knapsack-style child merges.
  \item Tree DP after bridge-tree or condensation reductions.
\end{itemize}

\section*{Practice Problems}

\begin{itemize}
  \item Maximum independent set on a tree.
  \item Count colorings with parent-child restrictions.
  \item Subtree optimization with one local state per node.
\end{itemize}

\section*{References}

\begin{itemize}
  \item \href{https://usaco.guide/gold/dp-trees}{USACO Guide: DP on Trees}.
  \item \href{https://codeforces.com/catalog}{Codeforces Catalog}.
\end{itemize}
